\section{Introduction}
\label{sec:intro}

Large language models are increasingly used as decision modules: a model reads a
case and emits an action that determines what a system does next. Under stochastic
decoding that action is a random variable, and repeated executions of the same input
can produce different actions
\citep{Ouyang_2025,atil2025nondeterminismdeterministicllmsettings}. What matters
operationally is not average accuracy over a benchmark but a joint property of the
individual decision: whether it is correct, and whether it recurs.

Consider an accounts-payable system in which a model decides whether an invoice
satisfies an approval policy. An invoice labelled \texttt{approve} on one execution
and \texttt{reject} on the next is not a reproducible decision. An invoice labelled
\texttt{reject} on all ten executions, wrongly, is not a usable one. In agentic
systems the decision selects a tool, a workflow branch, or the information passed
downstream, so both failure modes propagate. Operational reliability requires
correctness and reproducibility together, and aggregate accuracy reports neither.

Existing levers act on different parts of the stack. Batch-invariant kernels remove
numerical sources of run-to-run variation \citep{he2025nondeterminism}; fine-tuning
changes learned behaviour \citep{tian2023finetuninglanguagemodelsfactuality};
constrained and semantically guided decoding alter the generation procedure
\citep{loula2025syntactic,ahmed2025semanticprobabilisticcontrollanguage}; prompt
search reshapes the output distribution toward target tokens
\citep{bhargava2024whatsmagicwordcontrol}, and stability-guided refinement improves
consistency alongside task performance \citep{chen2025stability}. Each modifies
weights, numerics, the sampler, or the response after the fact. None answers what the
input must contain for a frozen model, under unmodified decoding and without retries,
to emit the intended action across independent samples. Sensitivity to phrasing is
well documented as a phenomenon
\citep{mizrahi-etal-2024-state,srikanth2024errorsnaturallanguagereasoning}; we treat
the task specification as the object of study rather than as a nuisance variable.
Earlier work named this \emph{interpretation drift} and showed that an explicit
specification brought four frontier models from $30\%$ agreement and 40 distinct
output hashes to a single hash across 40 runs on a classification task
\citep{nguyen2026substrates}.

We hold weights, decoding and retry policy fixed and vary only the specification
supplied before inference. The test case is a minimal operational decision: the
widely circulated car-wash prompt,\footnote{``I need to wash my car, the car wash is
only 50 metres away, should I walk or drive there?''} in which the destination is
within walking distance but completing the objective requires transporting the car,
so the intended action is \texttt{DRIVE}. Each model--condition cell is executed ten
times at temperature $1.0$ and classified by unanimity: \emph{reproducibly correct}
(the intended action in ten of ten samples), \emph{reproducibly incorrect} (zero of
ten), or \emph{non-reproducible} (anything between). We evaluate the untreated task,
seven conventional prompt interventions, and one shared task specification across 17
models.

No model is reproducibly correct on the untreated task: 13 of 17 are reproducibly
incorrect and four are non-reproducible. The seven conventional interventions raise
the aggregate rate from $7.6\%$ to at most $36.5\%$, and across all seven conditions
convert exactly one reproducibly incorrect model. The shared specification makes 14
of 17 reproducibly correct, $95.9\%$ of responses; the remaining three improve but do
not reach unanimity, so specification sufficiency is established per model--condition
rather than assumed. In open-weight models, transplanting a substrate-conditioned
activation into the untreated computation reverses its final
\texttt{DRIVE}--\texttt{WALK} preference; these traces use an earlier specification
than the fleet-wide experiment and are not claimed as the mechanism of the $14/17$
result.

\paragraph{Contributions.}
\begin{itemize}\itemsep2pt
\item An operational evaluation protocol that separates reproducibly correct,
reproducibly incorrect and non-reproducible decisions, and the finding that a single
shared task specification converts 11 reproducibly incorrect models to reproducibly
correct without weight updates, decoder changes, retries, majority voting or
post-generation correction.
\item Evidence that the structure of the supplied semantics is what matters:
instructions to reason step by step, adopt an expert role, avoid error, or apply
pressure convert one reproducibly incorrect model between them, while specifying the
task relations converts fourteen.
\item A causal measurement of the engineered transition: in open-weight models,
activation transplantation identifies internal states whose replacement reverses the
final action preference, without identifying the individual constraints or pathways
responsible.
\end{itemize}
